\documentclass[11pt,a4paper]{article}
\usepackage[utf8]{inputenc}
\usepackage[T1]{fontenc}
\usepackage[brazilian]{babel}
\usepackage[margin=2.3cm]{geometry}
\usepackage{xcolor}
\usepackage{hyperref}
\usepackage{enumitem}
\usepackage{booktabs}
\usepackage{tabularx}
\usepackage{longtable}
\usepackage{lmodern}
\usepackage{tcolorbox}
\usepackage{titlesec}
\usepackage{listings}

\hypersetup{colorlinks=true, urlcolor=blue!60!black, linkcolor=blue!50!black,
            citecolor=blue!50!black, pdfborder={0 0 0}}
\setlength{\emergencystretch}{3em}
\sloppy

\definecolor{accent}{HTML}{1F4E79}
\definecolor{ok}{HTML}{2E7D32}
\definecolor{warn}{HTML}{B45309}
\definecolor{codebg}{RGB}{245,247,250}
\definecolor{codeborder}{RGB}{210,215,225}
\definecolor{noticebg}{RGB}{232,242,250}
\definecolor{noticeborder}{RGB}{31,78,121}
\definecolor{warnbg}{RGB}{253,246,225}
\definecolor{warnborder}{RGB}{180,115,9}

\titleformat{\section}{\Large\bfseries\color{accent}}{\thesection}{0.7em}{}
\titleformat{\subsection}{\large\bfseries\color{accent!85!black}}{\thesubsection}{0.7em}{}

\newtcolorbox{notice}{
  colback=noticebg, colframe=noticeborder, boxrule=0.6pt,
  arc=2pt, top=4pt, bottom=4pt, left=8pt, right=8pt
}
\newtcolorbox{cuidado}{
  colback=warnbg, colframe=warnborder, boxrule=0.6pt,
  arc=2pt, top=4pt, bottom=4pt, left=8pt, right=8pt
}

\lstset{
  basicstyle=\ttfamily\scriptsize,
  backgroundcolor=\color{codebg},
  frame=single, framerule=0.4pt, rulecolor=\color{codeborder},
  framesep=4pt, xleftmargin=6pt, xrightmargin=4pt,
  columns=fullflexible, keepspaces=true, breaklines=true,
  literate=
    {á}{{\'a}}1 {é}{{\'e}}1 {í}{{\'i}}1 {ó}{{\'o}}1 {ú}{{\'u}}1
    {Á}{{\'A}}1 {É}{{\'E}}1 {Í}{{\'I}}1 {Ó}{{\'O}}1 {Ú}{{\'U}}1
    {â}{{\^a}}1 {ê}{{\^e}}1 {ô}{{\^o}}1
    {ã}{{\~a}}1 {õ}{{\~o}}1 {ç}{{\c{c}}}1 {Ç}{{\c{C}}}1,
}

\title{\textbf{Plano do ERP Patria --- v1}\\[0.4em]
  \large ERP grátis, multi-tenant SaaS e on-premise,
  monetizado via IA hipercomplexa de recomendação a fornecedores}
\author{Aarão Melo Lopes \and Claude (Patria Technology)}
\date{25 de maio de 2026}

\begin{document}
\maketitle

\begin{notice}
\textbf{Resumo executivo.} Construir um ERP completo (PDV, estoque, financeiro,
ordem de serviço, checkout) distribuído em dois modos a partir do mesmo
código: SaaS multi-tenant em \texttt{*.patriatechnology.com} e instalação
on-premise via Docker Compose na intranet do cliente. O produto é
\emph{grátis} para o lojista. Checkout entra 100\% funcional desde o dia 1
via \emph{conta interna} (Patria custodia o saldo em sua conta PagBank
principal e faz payout para a conta bancária do lojista; default D+1 com
opção de o lojista mudar para saque sob demanda) e oferece chave PagBank
própria como opt-in. Taxa por transação repassada ao lojista sem markup.
A receita vem da camada de IA: o NCO-HiperGNN (papers AF/AE/CT) ranqueia e
sugere produtos, e o fornecedor/distribuidor paga pela visibilidade
qualificada e por dados agregados anonimizados. Lançamento direto em
escala nacional, sem cold start regional.
\end{notice}

\tableofcontents
\newpage

%==================================================================
\section{Contexto e mercado-alvo}
%==================================================================

\subsection{Público-alvo}

Quatro segmentos, com prioridade comercial nesta ordem:

\begin{enumerate}[leftmargin=2em]
  \item \textbf{Comércio e materiais de construção} (lojas pequenas e médias).
        Mais carente de ERP decente, ticket médio razoável, baixa concorrência
        forte em N/NE.
  \item \textbf{Peças e serviços} (autopeças, oficina, eletrônica, marcenaria).
        Combina produto físico com mão de obra; OS é central.
  \item \textbf{Food service} (restaurantes, bares, lanchonetes), público da
        Pigz. Entrada posterior, com vantagem de já ter base de lojas instalada.
  \item \textbf{Eventos} (casas noturnas, festas, festivais), também mira
        adjacente à Pigz; ingressos e comanda eletrônica.
\end{enumerate}

\subsection{Diagnóstico do mercado}

O mercado de ERP brasileiro para PMEs é saturado em mensalidade fixa (Bling,
Omie, Tiny, Conta Azul, TOTVS de pequeno porte) e em PDV-via-maquininha
(Stone Ton, PagSeguro Moderninha, InfinitePay). Arquitetura, qualidade de UX
ou catálogo de features \emph{não} são diferenciadores defensáveis --- todos os
players conseguem construir tudo o que descrevemos como núcleo.

A brecha real está em \textbf{modelo de receita} e em \textbf{inteligência
preditiva matematicamente fundada}, não em features.

\subsection{Por que agora}

A Patria Technology já tem operando três peças que ninguém na concorrência
combina:

\begin{enumerate}[leftmargin=2em]
  \item Infraestrutura SaaS multi-tenant em produção (Patria + Patrícia),
        roteamento por subdomínio, worker headless via Claude Code,
        deploy automático.
  \item Solver neural NCO-HiperGNN treinado e validado (paper AF), com
        \emph{bound de tempo} ($\sim 10$ms) \emph{e bound de qualidade}
        ($C(s)^2/m^5$ rigoroso, paper AE) --- combinação que não existe
        em outro solver disponível no mercado.
  \item Integração PagBank em produção (PIX + cartão tokenizado), tabela
        \texttt{billing\_events}, webhook em
        \texttt{nco.patriatechnology.com}, validada e aprovada em 13/05/2026.
\end{enumerate}

%==================================================================
\section{Modelo de negócio}
%==================================================================

\subsection{Princípio}

\begin{notice}
\textbf{Tudo grátis para o lojista, inclusive financeiro.} Receita vem do
\textbf{lado do fornecedor}, mediada por inteligência preditiva.
\end{notice}

\subsection{Fontes de receita}

Cinco vetores, do mais simples ao mais complexo:

\begin{longtable}{@{}p{2.7cm}p{6.5cm}p{5.5cm}@{}}
\toprule
\textbf{Vetor} & \textbf{Como cobra} & \textbf{Pré-requisito} \\
\midrule
\endhead
Sugestão patrocinada de reposição & ``Tá acabando seu cimento --- sugiro
Votorantim com 7\% off agora.'' Fornecedor paga por sugestão aceita (CPA) ou
por lead qualificado (CPL). Sempre com \emph{flag patrocinado} visível ao
lojista. & Catálogo unificado de SKU + sinal de intenção de reposição (NCO
forecast). \\
\addlinespace
Marketplace de compras integrado & Lojista compra fornecedor dentro do ERP.
Take rate (1--3\%) sobre GMV, ou spread no crédito (parcelamento da compra).
& Conexão B2B com distribuidores; capital de giro para antecipar (já
disponível via fintech parceira). \\
\addlinespace
Insights agregados para indústria & Dashboard tipo Nielsen para Coca,
Brahma, Votorantim, indústrias de material e peças. ``Marca X cresceu 12\% no
varejo de bairro de Boa Vista nas últimas 4 semanas.'' Assinatura mensal
recorrente. & Massa crítica de lojas (alguns milhares), sync agregado
anonimizado opt-in. \\
\addlinespace
Trade marketing digital ao consumidor & Banner/cupom no cardápio digital,
totem, link de delivery do lojista. Marca paga CPM/CPC. & Canal de
consumidor final implementado (food/eventos). \\
\addlinespace
Crédito patrocinado pelo fornecedor & ``Coca te empresta R\$\,5k a juro
zero para repor estoque dela.'' Patria intermedeia, leva spread. & Score de
risco com base no histórico do ERP. \\
\bottomrule
\end{longtable}

\textbf{Fase 1 de receita:} foco em (1) e (2), os mais rápidos de validar.
(3) e (4) são consequência de massa instalada. (5) é prêmio depois da
maturidade.

\subsection{Por que isso é defensável}

\begin{itemize}[leftmargin=2em]
  \item \textbf{Versus ERPs tradicionais (Bling, Omie, Tiny, Conta Azul):}
        eles cobram mensalidade do lojista. Nós não cobramos. Quem está com
        caixa apertado nunca vai escolher pagar R\$\,200/mês se ``grátis''
        existe.
  \item \textbf{Versus fintechs com PDV (Stone, PagSeguro, InfinitePay):}
        elas vieram da maquininha; o ERP delas é fraco ou adendo. Nós viemos
        do ERP, que é o cérebro da loja. ERP é \emph{sticky}, maquininha é
        \emph{commodity}.
  \item \textbf{Versus Pigz:} a Pigz cobra mensalidade. Nós não. E nossa IA
        de recomendação tem fundação matemática rigorosa (bound de
        qualidade), não treino-e-funciona.
  \item \textbf{Versus qualquer recomendador clássico:} NCO-HiperGNN tem
        bound rigoroso de tempo \emph{e} qualidade. É o único quadrante
        marcado em \cite{paperAF}. Vira SLA contratual com fornecedor.
\end{itemize}

%==================================================================
\section{Diferenciais técnicos derivados do NCO}
%==================================================================

\subsection{O que o NCO-HiperGNN entrega ao ERP}

A camada de IA, executada localmente (CPU) ou via cloud (GPU em GEX44),
resolve quatro problemas industriais com latência $\sim 10$ms e variância
$7{,}5\times$ menor que heurísticas clássicas:

\begin{enumerate}[leftmargin=2em]
  \item \textbf{Forecast de venda por SKU.} Previsão de quantidade vendida
        no próximo período, derivada do histórico do ERP. Usa estrutura de
        cluster temporal $\to$ atribuição via partição.
  \item \textbf{Detecção de ruptura de estoque.} Probabilidade de
        zerar dada a curva de venda + lead time do fornecedor.
  \item \textbf{Segmentação automática de SKUs e clientes.} Grupos coesos
        para ações de venda (cross-sell, upsell), por simetria $S_3$ entre
        canais imaginários do quaternion.
  \item \textbf{Roteamento e escalonamento.} Coloração de grafo para
        agenda de entregas, OS, alocação de mesa (food).
\end{enumerate}

\subsection{Bound rigoroso vira SLA contratual}

Pela equação $C(s)^2 = 36(1-s^2)^2$ (paper AE) e arquitetura $m=4, L=4$, o
modelo entrega:

\begin{itemize}[leftmargin=2em]
  \item Latência $p_{99} < 50$ms para $n \le 500$ SKUs.
  \item Variância de output $\sigma_K \le C(s)/\sqrt{Km^5} \approx 0{,}07$
        para $K=8$ rollouts.
  \item Reprodutibilidade determinística com mesmo input + seed.
\end{itemize}

Esses três números viram cláusula de SLA para o cliente fornecedor que
patrocina sugestões.

\subsection{Decomposição micro/macro (paper CT)}

A teoria de campo médio descrita no paper CT permite servir dois lados
da plataforma de forma matematicamente coerente:

\begin{itemize}[leftmargin=2em]
  \item \textbf{Micro --- ao lojista.} Sugestão individual baseada na
        microdinâmica de cada elétron algébrico (paper A): ``compre Y,
        promova Z, repor X agora''.
  \item \textbf{Macro --- ao fornecedor.} Estado físico do agregado
        $S = \mathrm{mean}_k(s_k v_k)$ sobre toda a rede de lojas, com
        decomposição $S = s_\| V + S_\perp$. Defeito perpendicular
        $S_\perp$ revela \emph{heterogeneidade real} do mercado --- insight
        que Nielsen e GfK não entregam com essa granularidade.
\end{itemize}

São \emph{dois produtos sobre a mesma infraestrutura}, separáveis sem
perder coerência.

%==================================================================
\section{Arquitetura}
%==================================================================

\subsection{Topologia: um código, dois ambientes}

\begin{longtable}{@{}p{2.5cm}p{5.5cm}p{6.5cm}@{}}
\toprule
\textbf{Aspecto} & \textbf{SaaS} & \textbf{On-premise} \\
\midrule
\endhead
Hosting & Servidor Patria
(\texttt{erp.patriatechnology.com}, subdomínios) & Intranet do cliente
(\texttt{erp.empresa.local}) \\
Auth & Login web Patria + plano & Login interno do cliente \\
Multi-tenant & Sim, via header \texttt{X-Tenant} & Sim (uma instalação
serve N lojas da mesma empresa), \texttt{X-Tenant} interno; default
\texttt{tenant=default} se single-loja \\
Banco & Postgres central com pgvector & Postgres local no Docker Compose \\
LLM & Anthropic via Patria & BYOK (Anthropic/OpenAI key do cliente) ou Ollama
local opcional \\
NCO inference & Cloud (GEX44 ou réplica) & Container Python local com modelo
TorchScript step 80k empacotado (CPU-only acceptable) \\
Cloud sync & N/A & Opt-in. Sobe agregados anonimizados, baixa catálogo de
fornecedor e atualizações \\
Atualização & CI/CD automático & Manual via
\texttt{docker compose pull \&\& up -d} \\
\bottomrule
\end{longtable}

Diferença vive em variáveis de ambiente e em uma tabela \texttt{Installation}
com \texttt{mode}, \texttt{edition}, \texttt{cloudSyncEnabled}, \texttt{ncoMode},
\texttt{llmProvider}. Schema único.

\subsection{Estrutura de diretório}

\begin{lstlisting}
patria-erp/
|-- erp-api/                  NestJS + Prisma (já criado)
|-- erp-app/                  Vite + React + TS (já criado)
|-- erp-worker/               Prisma direto, Claude headless, jobs
|-- nco-inference/            Python + FastAPI, TorchScript step 80k
|-- catalog-cloud/            Catálogo unificado de fornecedor (apenas SaaS)
|-- docker-compose.yml        SaaS deployment
|-- docker-compose.local.yml  on-premise bundle
|-- install.sh                instalador on-premise
|-- estrategia/               este plano e futuros docs
\end{lstlisting}

\subsection{Princípios de design herdados do integrator Easysync}

A leitura do CLAUDE.md do Easysync integrator
(\texttt{/easysync/integrator/integrator/}) revela um ERP de varejo
maduro (Spring Boot, DB2, Liquibase, Kafka+Debezium) com domínio bem
modelado. Herdamos:

\begin{itemize}[leftmargin=2em]
  \item Divisão modular: \texttt{product, customer-supplier, stock,
        warehouse, service-order, budget, order, payment, cash, invoice,
        promotion, vehicle, warranty-term, notification, webhook}.
  \item Entidade \texttt{CustomerSupplier} unificada (cliente e fornecedor
        no mesmo cadastro --- decisão deles, mantida).
  \item Desenho do módulo \texttt{Cash}: \texttt{CashSession} com estados
        \texttt{OPEN/PARTIALLY\_CLOSED/CLOSED}, \texttt{CashOperation}
        para sangria/reforço/troca de operador, \texttt{CashPhysicalClose}
        por forma de pagamento.
  \item Filtros dinâmicos via \texttt{SearchFilter} (no NestJS, helper
        Prisma-where).
  \item Feature flags por módulo (\texttt{*\_ENABLED}).
\end{itemize}

Modernizamos:

\begin{longtable}{@{}p{5.5cm}p{8cm}@{}}
\toprule
\textbf{Integrator (referência)} & \textbf{ERP Patria (nosso)} \\
\midrule
\endhead
Spring Boot 3.3 / Java 21 & NestJS 11 / TypeScript / Node 22 \\
DB2 + Postgres secundário & Postgres único com pgvector \\
TF-IDF embeddings em \texttt{*\_EMBEDDING} & pgvector vector(768) + Gemini
embeddings, padrão Patria \\
Kafka + Debezium CDC & Outbox pattern interno + webhook HMAC, mais simples
para escala média \\
Chatbot TF-IDF + GPT & Patrícia (Claude + tools) + NCO-HiperGNN com bound
rigoroso \\
JWT manual com jjwt & \texttt{@nestjs/jwt} + Passport \\
iText (PDF) + Apache POI (Excel) + escpos-coffee & V2; MVP usa
puppeteer (PDF) + CSV; impressora térmica via driver web \\
GraalVM nativo opcional & NestJS standard (startup já rápido) \\
\bottomrule
\end{longtable}

\subsection{Stack tecnológica}

\begin{itemize}[leftmargin=2em]
  \item \textbf{Backend:} NestJS 11, TypeScript, Prisma 6, PostgreSQL 16 +
        pgvector, JWT via \texttt{@nestjs/jwt}, ConfigModule global.
  \item \textbf{Frontend:} Vite 5, React 19, TypeScript 5.6, proxy
        \texttt{/api $\to$ :3001}.
  \item \textbf{Worker:} Node 22, Prisma direto no banco, Claude Code
        headless via CLI (padrão patria-worker).
  \item \textbf{NCO inference:} Python 3.11, FastAPI, PyTorch (TorchScript
        congelado do checkpoint step 80k), CPU como default no on-premise.
  \item \textbf{Embeddings:} Gemini \texttt{gemini-embedding-001} 768d no
        SaaS; \texttt{paraphrase-multilingual-mpnet-base-v2} local no
        on-premise (mesmo padrão do
        \texttt{theory-embedder.service:9301} da Patria).
  \item \textbf{IA conversacional:} Claude Sonnet 4.6 (SaaS), Anthropic
        BYOK ou Ollama local (on-premise).
  \item \textbf{Pagamento:} integração PagBank reaproveitada da Patria
        (PIX + cartão tokenizado, tokenização client-side).
  \item \textbf{Deploy:} Docker Compose, imagens versionadas no registry.
\end{itemize}

%==================================================================
\section{Módulos de domínio}
%==================================================================

\subsection{Lista canônica}

\begin{longtable}{@{}p{3.5cm}p{6cm}p{4.5cm}@{}}
\toprule
\textbf{Módulo} & \textbf{Escopo} & \textbf{Status} \\
\midrule
\endhead
\texttt{auth} & login, JWT, roles, recuperação de senha & F1 \\
\texttt{installation} & singleton com modo, edição, feature flags & F1 \\
\texttt{tenant} & multi-tenant via \texttt{X-Tenant}, ALS interceptor & F1 \\
\texttt{customer-supplier} & cadastro unificado cliente/fornecedor & F1 \\
\texttt{product} & catálogo, hierarquia (Group/Subgroup/Section/Division),
embeddings pgvector & F1 \\
\texttt{warehouse} & depósitos & F1 \\
\texttt{stock} & saldo por SKU/depósito, movimento, ajuste & F1 \\
\texttt{promotion + price} & tabelas de preço, promoções & F2 \\
\texttt{order} & venda, item, desconto, parcela & F2 \\
\texttt{payment} & métodos de pagamento, conciliação & F2 \\
\texttt{checkout} & integração PagBank (PIX + cartão), tokenização
client-side, webhook & F2 \\
\texttt{cash} & sessão, sangria, reforço, troca de operador, fechamento
físico & F2 \\
\texttt{service-order} & OS, vinculação a veículo/equipamento, peça +
mão de obra & F3 \\
\texttt{vehicle} & veículo do cliente (oficina) & F3 \\
\texttt{warranty-term} & termos de garantia & F3 \\
\texttt{budget} & orçamento $\to$ efetiva em OS ou Order & F3 \\
\texttt{nco} & forecast, ruptura, cluster, sugestão & F4 \\
\texttt{catalog-sync} & sync opt-in cloud, ofertas patrocinadas & F4 \\
\texttt{invoice} & NF-e/NFC-e/NFS-e, integração SEFAZ & F5 \\
\texttt{webhook} & webhooks externos, HMAC, retry com backoff & F6 \\
\texttt{notification} & websocket, in-app, e-mail & F6 \\
\texttt{logistic} & frete, entrega & F6 \\
\bottomrule
\end{longtable}

\subsection{Modelo de dados resumido (núcleo)}

\begin{lstlisting}
Installation { id, mode (saas|on_premise), edition, flags, createdAt }
Tenant       { id, alias, name, status, plan, createdAt }
User         { id, tenantId, email, hash, role, occupationId }

CustomerSupplier { id, tenantId, type (CUSTOMER|SUPPLIER|BOTH),
                   name, document, contacts, address, embedding }

Product   { id, tenantId, sku, name, division, section, group, subgroup,
            unit, brand, gtin, costPrice, embedding, ... }
Warehouse { id, tenantId, name, address }
Stock     { id, tenantId, productId, warehouseId, qty, avgCost }
StockMovement { id, tenantId, productId, warehouseId, qty, type, refId }

Order        { id, tenantId, customerId, items, total, status,
               paymentId, cashSessionId }
OrderItem    { id, orderId, productId, qty, price, discount }
Payment      { id, tenantId, orderId, method, amount, status,
               pagbankOrderId, pagbankChargeId, paidAt,
               paymentMode (INTERNAL_WALLET|OWN_PAGBANK|OWN_OTHER) }
BillingEvent { id, tenantId, source, payload, processed, createdAt }

PaymentMode      { tenantId, mode, credentials (encrypted), createdAt }
MerchantWallet   { id, tenantId, currency, balance, blocked,
                   totalReceived, totalPaidOut }
WalletTransaction{ id, walletId, type, amount, balanceAfter,
                   refType, refId, createdAt }
PayoutPolicy     { tenantId, mode (AUTO_D1|MANUAL),
                   minAmount, bankAccountId }
BankAccount      { id, tenantId, type, bank, agency, account,
                   pixKey, pixKeyType, holderName, holderDocument }
Payout           { id, tenantId, walletId, amount, fee, netAmount,
                   status, scheduledFor, processedAt,
                   externalTxId, failureReason }

Cash             { id, tenantId, name }
CashSession      { id, cashId, openedBy, openedAt, closedAt, status }
CashOperation    { id, sessionId, type, amount, operatorId, createdAt }
CashPhysicalClose{ id, sessionId, method, amount, diff }

ServiceOrder { id, tenantId, customerId, vehicleId, items[],
               labor[], status, warrantyTermId }
Vehicle      { id, tenantId, customerId, plate, model, brand, year }
Budget       { id, tenantId, customerId, items, total, status,
               convertedToId, convertedToType }

NcoJob       { id, tenantId, type, input, output, latencyMs,
               modelVersion, createdAt }
CatalogSyncLog { id, direction, kind, count, createdAt }
\end{lstlisting}

%==================================================================
\section{Checkout: conta interna + chave própria opcional}
%==================================================================

\subsection{Princípio: fricção zero}

\begin{notice}
\textbf{Filosofia.} O lojista entra e vende no minuto 1, sem precisar
abrir conta PagBank, sem documentação, sem espera de homologação. O ERP
usa a conta PagBank principal da Patria como \emph{custodiante técnico} e
mantém um \emph{ledger interno} por lojista. Payouts automáticos
transferem o saldo do lojista para a conta bancária dele cadastrada no
onboarding (D+1 default; lojista pode optar por acumular). Lojista que
prefira autonomia configura a chave PagBank própria a qualquer momento e
o ERP passa a operar em BYOK. Taxa por transação é repassada ao lojista
sem markup --- a Patria não lucra na taxa.
\end{notice}

\subsection{Base técnica reaproveitada da Patria}

A integração PagBank foi homologada em produção em 13/05/2026 (anexo v4,
\texttt{nco.patriatechnology.com}). Já cobre:

\begin{itemize}[leftmargin=2em]
  \item \textbf{PIX:} \texttt{POST /api/billing/checkout} cria order com
        \texttt{qr\_codes}; PagBank retorna copia-cola + PNG.
  \item \textbf{Cartão de crédito:} tokenização client-side via SDK
        \texttt{pagseguro.min.js}; backend recebe apenas string
        \texttt{encrypted}; PAN e CVV nunca trafegam pelo backend.
  \item \textbf{Webhooks:} \texttt{POST /api/webhooks/pagbank} grava em
        \texttt{billing\_events} (\texttt{processed=false}), depois processa
        e marca como \texttt{processed=true}. Idempotência via
        \texttt{reference\_id = pat-<tenant>-<timestamp>}.
  \item \textbf{Conta principal:} \texttt{patriatechnology53@gmail.com}
        (PROD).
\end{itemize}

Essa stack é extraída para módulo compartilhado e usada pelo ERP em dois
cenários: cobrança da mensalidade Patria (quando houver edição paga) e
cobrança do cliente final do lojista no PDV/OS.

\subsection{Três modos de operação por tenant}

\begin{longtable}{@{}p{3.5cm}p{5.5cm}p{5.5cm}@{}}
\toprule
\textbf{Modo} & \textbf{Como funciona} & \textbf{Quando usar} \\
\midrule
\endhead
\texttt{INTERNAL\_WALLET} (default) & Recebimento cai na conta PagBank
Patria; ledger interno credita o lojista; payout D+1 (default) ou sob
demanda para a conta bancária dele. & Lojista pequeno/médio que quer
começar a vender hoje sem burocracia. \\
\addlinespace
\texttt{OWN\_PAGBANK} (BYOK) & Lojista configura token PagBank próprio
no onboarding ou nas configurações; checkout opera diretamente na conta
dele. & Lojista que já tem conta PagBank, prefere autonomia, ou quer
evitar custódia. \\
\addlinespace
\texttt{OWN\_OTHER} (futuro) & Slot para outras integrações: Mercado Pago,
Stone, Asaas, InfinitePay. Mesma interface, adapter por provedor. & Quando
houver demanda real. \\
\bottomrule
\end{longtable}

A configuração vive em \texttt{PaymentMode} por tenant; mudar de modo é
operação manual com validação (não pode mudar com saldo virtual maior
que zero sem saque antes).

\subsection{Ledger interno (modo \texttt{INTERNAL\_WALLET})}

A custódia técnica do saldo na conta PagBank Patria é refletida em um
ledger contábil de dupla entrada simplificado:

\begin{lstlisting}
MerchantWallet {
  id, tenantId, currency='BRL',
  balance,         -- saldo disponivel para payout
  blocked,         -- reservado (pre-autorizacao, contestacao pendente)
  totalReceived,   -- acumulado vitalicio recebido
  totalPaidOut,    -- acumulado vitalicio pago ao lojista
  createdAt, updatedAt
}

WalletTransaction {
  id, walletId, type, amount, balanceAfter,
  -- types:
  --   SALE_CREDIT       (venda confirmada via webhook PAID)
  --   FEE_DEBIT         (taxa PagBank repassada)
  --   PAYOUT_DEBIT      (transferencia ao lojista)
  --   PAYOUT_FEE_DEBIT  (taxa do payout, se Pix tiver custo)
  --   REFUND_DEBIT      (estorno parcial/total)
  --   CHARGEBACK_DEBIT  (contestacao ganha pelo cliente final)
  --   MANUAL_ADJUST     (correcao operacional, motivo obrigatorio)
  refType, refId,  -- ligacao ao Order/Payment/Payout originador
  createdAt
}
\end{lstlisting}

\subsection{Política de payout (modelo híbrido)}

\begin{itemize}[leftmargin=2em]
  \item \textbf{Default D+1.} Toda venda confirmada hoje vira payout
        automático no próximo dia útil para a conta bancária cadastrada.
        Minimiza tempo de custódia (e portanto exposição da Patria) e dá
        previsibilidade ao lojista.
  \item \textbf{Opt-in: acumular.} Lojista pode mudar para modo manual,
        em que o saldo acumula e ele clica em ``sacar'' quando quiser
        (mínimo configurável, ex.\ R\$\,50).
  \item \textbf{Worker} \texttt{erp-worker} roda job diário 06:00:
        consolida saldo, gera \texttt{Payout} pendente, dispara Pix da
        conta Patria para a conta do lojista, marca \texttt{PROCESSED} no
        retorno.
\end{itemize}

\begin{lstlisting}
PayoutPolicy {
  tenantId,
  mode (AUTO_D1 | MANUAL),
  minAmount,            -- so dispara payout >= esse valor
  bankAccountId         -- destino
}

BankAccount {
  id, tenantId, type (CHECKING|SAVINGS|PIX),
  bank, agency, account, accountDigit,
  pixKey, pixKeyType (CPF|CNPJ|EMAIL|PHONE|RANDOM),
  holderName, holderDocument
}

Payout {
  id, tenantId, walletId, amount, fee, netAmount,
  status (PENDING | PROCESSING | PAID | FAILED | CANCELLED),
  scheduledFor, processedAt,
  externalTxId,         -- end-to-end ID do Pix executado
  failureReason
}
\end{lstlisting}

\subsection{Taxa: repasse sem markup}

A taxa cobrada pelo PagBank por transação (PIX \textasciitilde 0{,}99\%,
crédito 1x \textasciitilde 3{,}19\%, com variações por modalidade) é
\textbf{repassada integralmente ao lojista, sem markup}. Patria não lucra
na taxa --- isso preserva a promessa ``grátis'' (no sentido honesto: o
lojista paga exatamente o que pagaria se tivesse conta PagBank própria).
A receita Patria continua vindo da camada de IA / fornecedor, não de
spread financeiro.

A \texttt{FEE\_DEBIT} é lançada na mesma transação do \texttt{SALE\_CREDIT}
e o lojista vê o valor líquido na tela.

\subsection{Fluxos suportados no PDV}

\begin{itemize}[leftmargin=2em]
  \item \textbf{PIX no balcão:} backend gera QR code, balcão exibe (TV,
        totem ou impressora térmica), cliente paga, webhook confirma,
        \texttt{Payment.status=PAID}, baixa de estoque, recibo, e
        \texttt{SALE\_CREDIT} no \texttt{MerchantWallet} (modo interno) ou
        crédito direto na conta PagBank do lojista (BYOK).
  \item \textbf{Cartão presencial:} maquininha física PagBank (Order API).
        Status sincrônico + webhook redundante.
  \item \textbf{Cartão remoto / link de pagamento:} URL gerada pelo
        backend, cliente paga sem estar fisicamente na loja (forte para
        material de construção, oficina e venda B2B).
  \item \textbf{Crediário interno:} parcela registrada em
        \texttt{AccountsReceivable}, sem trânsito por PagBank
        (instrumento próprio do lojista, sem custódia).
\end{itemize}

\subsection{Considerações on-premise}

\begin{itemize}[leftmargin=2em]
  \item Cliente on-premise pode operar em modo \texttt{OWN\_PAGBANK}
        configurando \texttt{PAGBANK\_TOKEN} no \texttt{.env} --- checkout
        idêntico ao SaaS, sem dependência da Patria central.
  \item Cliente on-premise pode operar em \texttt{INTERNAL\_WALLET} se
        ativar \texttt{CLOUD\_SYNC\_ENABLED=true} e aceitar termo
        específico de custódia --- requer reconciliação periódica com o
        cloud Patria.
  \item Sem internet, checkout PagBank fica offline; crediário e
        dinheiro continuam funcionais.
\end{itemize}

\subsection{Migração futura para split nativo (PagBank Plus/Connect)}

Quando o volume justificar e o PagBank Plus/Connect (ou equivalente)
estiver integrado, o modo \texttt{INTERNAL\_WALLET} migra para
\emph{split nativo}: o pagamento já é dividido na origem entre a conta
do lojista (subaccount no PagBank/BaaS) e uma eventual conta de taxa da
Patria. Custódia some, peso regulatório some, ledger continua existindo
mas vira espelho do que o PagBank/BaaS reporta. Migração transparente
para o lojista; nenhuma mudança de UX.

%==================================================================
\section{Sync cloud opt-in}
%==================================================================

\subsection{Motivação}

Sem dado agregado, vetores de receita (3) e (4) não funcionam, e a IA
não melhora. Mas o on-premise por contrato \emph{não} pode forçar
nada. Solução: opt-in transparente.

\subsection{Modelo do que sobe}

Apenas \textbf{agregados anonimizados}. Endpoint
\texttt{GET /admin/cloud-sync/preview} mostra ao cliente exatamente o que
seria enviado antes de habilitar.

\begin{itemize}[leftmargin=2em]
  \item Vendas agregadas por SKU/categoria/semana (sem PII, sem valor
        individual).
  \item Eventos de ruptura.
  \item Sinais de intenção de reposição.
  \item Métricas operacionais (tempo médio de OS, taxa de cancelamento).
\end{itemize}

\subsection{Modelo do que desce}

\begin{itemize}[leftmargin=2em]
  \item Catálogo enriquecido de fornecedor (códigos, fotos, descrições,
        equivalências, ofertas patrocinadas).
  \item Atualizações de modelo NCO (novos checkpoints).
  \item Patches de aplicação (opt-in).
\end{itemize}

\subsection{Privacidade}

\begin{cuidado}
\textbf{LGPD.} Nenhum dado pessoal de cliente final do lojista
(CPF, nome, telefone, endereço) sobe ao cloud em nenhum cenário. Apenas
agregados estatísticos. O termo de uso explicita o que é coletado,
para qual finalidade, e o cliente pode revogar a qualquer momento.
Acordo de processamento de dados (DPA) padrão.
\end{cuidado}

%==================================================================
\section{Fases de entrega}
%==================================================================

\subsection{F1 --- Núcleo comum (auth + cadastros + estoque)}

\textbf{Escopo.} \texttt{auth, installation, tenant, customer-supplier,
product, warehouse, stock}.

\textbf{Entregáveis.}
\begin{itemize}[leftmargin=2em]
  \item Login/logout, recuperação de senha.
  \item Cadastro de empresa (Installation) e usuário inicial admin.
  \item CRUD de cliente/fornecedor unificado.
  \item CRUD de produto com hierarquia Division/Section/Group/Subgroup.
  \item Embedding pgvector na criação/atualização de produto e
        customer-supplier (Gemini no SaaS, modelo local no on-premise).
  \item CRUD de depósito.
  \item Saldo de estoque por SKU/depósito + movimentação manual (entrada,
        saída, ajuste).
\end{itemize}

\textbf{Critério de aceite.}
\begin{itemize}[leftmargin=2em]
  \item Multi-tenant funcional (header \texttt{X-Tenant} isola dados em
        \emph{todos} os endpoints).
  \item Embeddings populados; \texttt{GET /products/search?q=...} retorna
        por similaridade.
  \item Build SaaS e build on-premise (Docker Compose) ambos rodando.
\end{itemize}

\subsection{F2 --- Vendas + Checkout + Caixa}

\textbf{Escopo.} \texttt{order, payment, checkout, cash, promotion, price}.

\textbf{Entregáveis.}
\begin{itemize}[leftmargin=2em]
  \item Venda com múltiplos itens, descontos por item e total.
  \item Tabela de preços + promoções por período.
  \item Caixa: \texttt{CashSession} (open/partially\_closed/closed),
        sangria, reforço, troca de operador (atômico), fechamento físico
        por método.
  \item \textbf{Checkout 100\% funcional} com os três modos:
        \texttt{INTERNAL\_WALLET} (default), \texttt{OWN\_PAGBANK} (BYOK),
        \texttt{OWN\_OTHER} (slot para futuro). PIX balcão, cartão
        presencial via Order API, link de pagamento remoto.
  \item \textbf{Wallet interno}: \texttt{MerchantWallet} +
        \texttt{WalletTransaction} + reconciliação automática com webhook
        PagBank.
  \item \textbf{Payout automático D+1} via worker (Pix da conta Patria
        para a conta bancária do lojista); opt-in de modo manual ``sacar
        quando quiser''.
  \item Onboarding: lojista cadastra \texttt{BankAccount}, escolhe modo
        (interno é default), aceita termo de custódia se for interno.
  \item Repasse de taxa PagBank ao lojista sem markup, visível como
        \texttt{FEE\_DEBIT} na timeline da wallet.
  \item Crediário interno (\texttt{AccountsReceivable} básico, parcelas).
  \item PDV web responsivo no \texttt{erp-app}.
\end{itemize}

\textbf{Critério de aceite.}
\begin{itemize}[leftmargin=2em]
  \item Venda completa: scan/SKU $\to$ carrinho $\to$ pagamento
        (PIX/cartão/dinheiro) $\to$ baixa de estoque $\to$ recibo $\to$
        crédito no wallet do lojista (modo interno) ou na conta dele
        (BYOK).
  \item Webhook PagBank atualiza \texttt{Payment.status},
        \texttt{WalletTransaction} \texttt{SALE\_CREDIT}+\texttt{FEE\_DEBIT}
        atômicos.
  \item Job de payout D+1 executa, faz Pix, lança \texttt{PAYOUT\_DEBIT};
        lojista vê dinheiro na conta bancária no dia seguinte.
  \item Tenant pode alternar entre modos (com restrição: saldo zerado).
  \item Sessão de caixa abre, opera, fecha; relatório de fechamento bate
        contra contagem física.
\end{itemize}

\subsection{F3 --- Vertical OS (peças + serviços)}

\textbf{Escopo.} \texttt{service-order, vehicle, warranty-term, budget}.

\textbf{Entregáveis.}
\begin{itemize}[leftmargin=2em]
  \item Cadastro de veículo/equipamento por cliente, histórico de OS.
  \item OS com peça + mão de obra no mesmo documento, status
        (aberta/em andamento/aguardando peça/finalizada/entregue).
  \item Orçamento que converte em OS ou Order.
  \item Termo de garantia anexável à OS.
\end{itemize}

\textbf{Critério de aceite.}
\begin{itemize}[leftmargin=2em]
  \item Ciclo orçamento $\to$ aprovação $\to$ execução $\to$ entrega
        funcional para oficina/autopeças piloto.
  \item Cálculo de baixa de estoque correto na finalização.
  \item PDF de OS gerado e impresso.
\end{itemize}

\subsection{F4 --- IA: NCO + sync de catálogo}

\textbf{Escopo.} \texttt{nco, catalog-sync}.

\textbf{Entregáveis.}
\begin{itemize}[leftmargin=2em]
  \item Container \texttt{nco-inference} empacotado, modelo TorchScript
        step 80k embutido, endpoints \texttt{POST /partition},
        \texttt{POST /forecast}, \texttt{POST /cluster}.
  \item Worker \texttt{erp-worker} agendado: forecast noturno de top-N
        SKUs, detecção de ruptura, geração de \texttt{NcoJob}.
  \item Sugestão patrocinada de reposição na tela do PDV (com flag visual
        ``patrocinado'').
  \item \texttt{catalog-sync} bidirecional, opt-in, com endpoint de preview.
\end{itemize}

\textbf{Critério de aceite.}
\begin{itemize}[leftmargin=2em]
  \item Previsão de ruptura disparada com lead time suficiente para
        reposição em pelo menos 80\% dos SKUs de giro alto da loja piloto.
  \item Sugestão de fornecedor aparece no carrinho com ranking explicado.
  \item Sync funciona ida e volta em loja com internet intermitente.
\end{itemize}

\subsection{F5 --- Fiscal (NF-e/NFC-e/NFS-e)}

\textbf{Escopo.} \texttt{invoice}.

\textbf{Entregáveis.}
\begin{itemize}[leftmargin=2em]
  \item Integração com homologação SEFAZ por estado (UF do tenant).
  \item Emissão de NF-e (B2B), NFC-e (B2C presencial), NFS-e (serviço,
        para OS).
  \item Cancelamento, carta de correção, contingência.
\end{itemize}

\textbf{Crítico.} Esta fase tem dependência regulatória pesada. Usar
biblioteca consolidada (nfe.io, FocusNFe, eNotas) como primeira opção;
implementação própria só se justificar economicamente.

\subsection{F6 --- Extras}

\textbf{Escopo.} \texttt{webhook, notification, logistic, integration}.

\begin{itemize}[leftmargin=2em]
  \item Webhooks externos para o cliente (loja envia evento para sistema
        próprio).
  \item Notificações em tempo real (websocket), in-app, e-mail.
  \item Frete (cálculo via integração com transportadora ou Correios) e
        rastreio.
\end{itemize}

%==================================================================
\section{Riscos e mitigações}
%==================================================================

\begin{longtable}{@{}p{4.5cm}p{4.5cm}p{5.5cm}@{}}
\toprule
\textbf{Risco} & \textbf{Probabilidade/impacto} & \textbf{Mitigação} \\
\midrule
\endhead
Adoção lenta por desconfiança (``ERP grátis é golpe'') & Alta / alto &
Reputação Patria, casos piloto com depoimento, transparência total no
modelo de receita, código aberto opcional. \\
\addlinespace
\textbf{Regulatório BACEN} no modo \texttt{INTERNAL\_WALLET} (Patria
custodia saldo de terceiros na conta PagBank principal) & Média / muito
alto &
Quatro mitigações combinadas:
(1) \textbf{Payout D+1 agressivo} mantém tempo médio de custódia
$< 24$h, minimizando saldo agregado retido.
(2) \textbf{Limite por tenant} configurável (ex.\ R\$\,50k de saldo
máximo); acima disso forçar BYOK ou saque imediato.
(3) \textbf{Consulta a advogado especialista em fintech} antes de
operar acima de saldo agregado de R\$\,1\,M ou
$> 100$ tenants ativos no modo interno.
(4) \textbf{Plano explícito de migração} para PagBank Plus/Connect
ou BaaS (Asaas/Pagar.me/Celcoin) com sub-account emitida por IP
autorizado, eliminando custódia. \\
\addlinespace
Fornecedor não topa pagar até massa crítica & Alta / alto &
Receita inicial não depende disso. Lojista já recebe valor real (forecast,
ruptura, sugestão sem patrocínio) desde dia 1. Receita do fornecedor é
camada 2, aceitar latência. \\
\addlinespace
Cliente percebe ranking como ``vendido'' & Média / alto &
Flag \emph{patrocinado} sempre visível, ranking transparente (qualidade
empírica × lance, como Google Ads), publicar metodologia. \\
\addlinespace
LGPD: vazamento ou processamento indevido & Baixa / muito alto &
Anonimização rigorosa no sync, opt-in explícito, DPA padrão, criptografia
em trânsito e em repouso, log de acesso. \\
\addlinespace
Suporte on-premise difícil (debug remoto) & Alta / médio &
Telemetria opt-in (apenas erros, sem dado), comando \texttt{erp-doctor}
gera bundle de diagnóstico anônimo para enviar ao suporte. \\
\addlinespace
NCO-HiperGNN out-of-distribution colapsa em domínio novo & Média / alto &
Treinar na distribuição do cliente (custo $\sim 3$ dias GPU), conforme
limitação documentada no paper AF; modelo fallback (heurístico) para
domínios sem treino. \\
\addlinespace
Concorrência (Stone, PagSeguro) replica modelo & Média / alto &
Defensável por: bound matemático do NCO (paper AE), Patrícia
(orquestração via Claude), tempo de mercado. Velocidade > perfeição. \\
\addlinespace
Bug de fiscal causa NF-e rejeitada em produção & Alta / alto &
Usar provedor consolidado (FocusNFe etc.) até maturidade, fila de
retentativa, alerta para diretor responsável. \\
\bottomrule
\end{longtable}

%==================================================================
\section{Próximos passos imediatos}
%==================================================================

\begin{enumerate}[leftmargin=2em]
  \item \textbf{Validar este plano.} Aarão revisa, ajusta escopo das fases,
        confirma ordem de prioridade.
  \item \textbf{Definir tenant inicial de teste.} Uma loja real (idealmente
        material de construção ou autopeças em Boa Vista) que aceite
        colaborar como piloto de F1+F2.
  \item \textbf{Começar F1.} Estruturar módulos \texttt{auth, installation,
        tenant, customer-supplier, product, warehouse, stock} no erp-api
        já criado. Schema Prisma definitivo. Primeira tela do erp-app.
  \item \textbf{Preparar nco-inference.} Exportar TorchScript do checkpoint
        step 80k do GEX44, validar inferência local CPU.
  \item \textbf{Reuso PagBank.} Extrair módulo \texttt{checkout} da Patria
        para biblioteca compartilhada (\texttt{@patria/checkout}?) ou
        copiar verbatim para erp-api. Adicionar camada
        \texttt{wallet} + \texttt{payout} acima da integração crua.
  \item \textbf{Termo de custódia (modo interno).} Redigir termo
        específico explicando que a Patria custodia o saldo na conta
        PagBank principal e faz payout D+1, sem juros, sem retenção
        operacional, sem cobrança extra. Aceite obrigatório no onboarding
        de tenant que escolher modo \texttt{INTERNAL\_WALLET}. Consulta
        rápida com advogado fintech antes de publicar.
\end{enumerate}

%==================================================================
\section*{Referências internas}
%==================================================================

\begin{thebibliography}{9}
\bibitem{paperAF} Lopes, A.; Claude.
\emph{Paper AF --- NCO-HiperGNN: solver neural para problemas combinatoriais
de partição via álgebra hipercomplexa generalizada}, 2026.
\bibitem{paperAE} Lopes, A.; Claude.
\emph{Paper AE --- Ponte Schur-Weyl: Hochschild $\leftrightarrow$ tensorial},
2026.
\bibitem{paperCT} Lopes, A.; Claude.
\emph{Paper CT --- Mecânica de campo médio em álgebras hipercomplexas},
2026.
\bibitem{ncoPartition} Lopes, A.; Claude.
\texttt{NCO\_PARTITIONING.md} --- roadmap do produto NCO, 2026.
\bibitem{pagbankV4} Patria Technology.
\emph{Anexo v4 --- Validação PagBank em PRODUÇÃO}, 13/05/2026.
\bibitem{integratorClaude} Easysync.
\texttt{integrator/CLAUDE.md} --- referência de arquitetura ERP de varejo
maduro.
\bibitem{patriaProduct} Patria Technology.
\texttt{patria-api/PRODUCT.md} --- visão do produto Patria.
\end{thebibliography}

\end{document}
